Why is it not enough to see the flower alone? This paper takes that question seriously and develops, in answer to it, a systematic reconstruction of eudaimonia under the framework of Generative Relational Being (GRB). We argue that both the existentialist tradition---from Heidegger's \emph{Mitsein} to Buber's \emph{I--Thou}---and the eudaimonist tradition---from Aristotle's \emph{phronesis} to Csikszentmihalyi's flow and Seligman's PERMA---share a structural presupposition that prevents them from accounting for shared happiness: the assumption that the subject of flourishing is the individual.

Against this presupposition, we advance three interlocking claims. First, \emb{the subject of contingency is the relation}: the flower descends not upon either person but upon the relational field constituted by their shared life, and the act of sharing is the constitutive act by which a contingent event becomes a relational event. Second, \emb{sharing is not the expression of relational happiness but its generative mechanism}: it triggers co-evolutionary dynamics in the coupled relational system that produce a shared attractor irreducible to either individual's dynamics. Third, \emb{relational happiness is an emergent signal of co-evolutionary activity}: it cannot be possessed by either party because it does not reside in either party.

These claims are grounded in coupled dynamical systems theory (Kuramoto synchronisation, quantum entanglement structure), neuroscience (EEG/MEG/fMRI/fNIRS hyperscanning, embodied resonance, co-regulation), and a formal analogy with spiking neural networks (STDP as the mechanism of relational structural modification). We propose an empirical research programme for testing these claims, undertake a hermeneutic re-reading of eleven classical happiness concepts under GRB, and situate the account cross-culturally through Japanese \emph{ma} and \emph{en}, Confucian \emph{ren}, and Ubuntu philosophy. The paper closes with a praxis chapter on the daily cultivation of relational happiness and a literary envoi that returns to the flower.
